\section{Formalization, the use of AI, and open questions}\label{sec:formalization}

\subsection{The formalization}

The repository~\cite{Companion} (Apache-2.0) contains five Lean~4 libraries on
top of Mathlib (Lean and Mathlib \texttt{v4.35.0-rc3}):
\begin{itemize}
\item \texttt{MovingSofaOptimality} (51 files, about 40,000 lines): Baek's
paper~\cite{Baek}, organized by its chapters, with the results that it cites
(Schneider's area formula for planar convex bodies; the existence and uniqueness
of the solution of Romik's system in the box of \cref{fact:gerver}) and the
structure of Gerver's sofa
(\cref{fact:gerver}\ref{fact:gerver:cap},\ref{fact:gerver:niche}), which
\cite{Baek} states without proof (Thm.~8.4.1) or proves invalidly (Thm.~6.1.2);
\item \texttt{MovingSofaUniqueness} (11 files, about 6,000 lines):
\cref{sec:selection,sec:variations,sec:right-angle,sec:injectivity,sec:equality,sec:recovery},
one module per step of the argument, the module \texttt{Maximizing}, shared by
\cref{rem:second} and \cref{sec:unified} (the maximizing caps, and the steps
from their two properties in \cref{prop:unified-caps} to optimality and
uniqueness), and the module \texttt{MaximizerRoute} (the rest of the second
proof of optimality);
\item \texttt{MovingSofaBridge} (4 files, about 2,600 lines): the connection
with the definitions of formal-conjectures~\cite{FormalConjectures}, which are
different (a motion there is a path of isometries that starts at the identity,
and Gerver's sofa is defined from Gerver's four constants); it shows that a set
is a moving sofa for formal-conjectures if and only if it lies in the horizontal
side of the hallway and is a moving sofa in the sense of
Appendix~\ref{app:lean}, and that the optimal areas and the Gerver's sofas of
the two formulations agree, so that formal-conjectures' statements follow from
Baek's theorem and from \cref{thm:main}, with the existence and uniqueness of
Gerver's four constants, which the bridge proves; the uniqueness statement,
which it lists as open, follows from \cref{thm:main};
\item \texttt{MovingSofaStability} (12 files, about 7,600 lines):
\cref{sec:stability}, one module per step of the proof, and the certificate of
\cref{sec:unified} (\cref{thm:certificate});
\item \texttt{MovingSofaExtremal} (3 files, about 240 lines):
\cref{sec:unified}, with the steps that it shares with the second proof of
optimality imported from \texttt{MovingSofaUniqueness};
\texttt{MovingSofaStability} takes its optimality and uniqueness theorems from
\texttt{MovingSofaExtremal}.
\end{itemize}
The file \texttt{Challenge.lean} states fifteen theorems, among them Baek's
theorem (\path{Baek.gerver_sofa_optimal}), \cref{thm:main}
(\path{Baek.gerver_sofa_unique}), \cref{thm:stability},
\cref{thm:stability-angle} and
\cref{thm:stability-exponent}\ref{thm:stability-exponent:no}
(\path{Baek.gerver_sofa_stable}, \path{Baek.gerver_sofa_angle_stable},
\path{Baek.gerver_sofa_stability_exponent}). The eight theorems of the namespace
\texttt{Baek} use only Mathlib's vocabulary and the definitions of
Appendices~\ref{app:gerver} and~\ref{app:lean}; the other seven use the
definitions of formal-conjectures, restated verbatim in a namespace of their own
(an anonymous topology instance is given a name), and the three theorems of the
namespace \texttt{Bridge} also use Baek's definitions. \texttt{Solution.lean},
with the modules that it imports, proves them: the Challenge is a file of
statements, so its fifteen theorems end in \texttt{sorry} by design, and
Comparator checks that \texttt{Solution.lean} proves exactly the statements of
\texttt{Challenge.lean}. \texttt{SolutionCoercive.lean} proves the fifteen
statements again, in a namespace of its own, taking optimality, uniqueness and
the stability of \cref{thm:stability,thm:stability-angle} from
\texttt{MovingSofaExtremal}; it imports neither \texttt{Solution.lean} nor the
module \texttt{Main} of \texttt{MovingSofaUniqueness}. \texttt{lake build}
succeeds, and an audit checks that every declaration of the five libraries and
of \texttt{SolutionCoercive.lean} depends only on the axioms
\texttt{propext}, \texttt{Classical.choice} and \texttt{Quot.sound}. The project
is registered in the Palomar registry as entry PALOMAR-2026-10-02-000008. Its
version~4 registers an earlier commit, whose Challenge states the twelve
theorems other than the three of stability. The continuous integration passed on
the commit cited in~\cite{Companion}. To rebuild, run \texttt{lake exe cache get
\&\& lake build}, then the commands listed in the repository's
\texttt{README.md}. The Challenge states \cref{thm:main}, \cref{thm:stability},
\cref{thm:stability-angle} and
\cref{thm:stability-exponent}\ref{thm:stability-exponent:no}; the other results
of this paper are proved in the libraries, and Appendix~\ref{app:lean} lists the
declarations.

The machine check shows that \cref{thm:main}, the theorems of
\cref{sec:stability} and those of \cref{sec:unified}, as the Lean statements
formulate them, follow from Lean's axioms and Mathlib. The formal proof
contains Baek's argument and the argument of this paper. The machine check does
not check that the definitions say what
they should. Appendix~\ref{app:lean} reproduces the Lean definitions (moving
sofa, hallway, the rotation path, Gerver's sofa, and the definitions used by the
stability theorems) in mathematical notation, so that they can be read and
compared with the intended meaning. Every proof of a result of~\cite{Baek} in
the library follows Baek's argument, except at sixteen steps, each forced by an
error or gap of the paper, by mathematics that Mathlib lacks (the Jordan curve
theorem, Green's theorem, the Brunn--Minkowski inequality, mixed volumes), or by
the representations chosen in the formalization (the surface area measure as a
Lebesgue--Stieltjes measure, the area functional of a convex arc as $\frac12\int
h\dd\sigma$, no orientations of curves). A route check in continuous integration
compares the results used by each Lean proof with those cited by Baek's proof;
the other details were compared with Baek's TeX source by AI agents, not by a
person.

The second proof of optimality (\cref{rem:second}) is formalized in two modules
of \texttt{MovingSofaUniqueness}. The module \texttt{MaximizerRoute} proves the
two properties of the maximizing caps in \cref{prop:unified-caps} with Baek's
bound and \cref{prop:kernel} in place of the certificate. The module
\texttt{Maximizing}, which \texttt{MovingSofaExtremal} uses too, derives from
these two properties the assertions of \cref{thm:unified-optimality} and, again,
\cref{thm:main} and \cref{cor:all}. Neither imports the module
\texttt{Main} of \texttt{MovingSofaUniqueness}, which proves \cref{thm:main}
from Baek's theorem.
A separate audit checks that none of the 35 declarations of these two modules
depends, directly or through other declarations, on \texttt{Main}, on
\baek{Thm.~1.1.1}, or on the results from which Baek derives step~(3) for the
caps of step~(2):
\baek{Thm.~1.5.2, Thm.~8.1.1\,(2)} and, among the results they rest on,
\baek{Thm.~4.1.2, 4.1.4, 4.2.5, 6.1.1, 6.3.3, 6.4.3, 6.5.6} and
\baek{Cor.~6.4.4}.

The certificate of \cref{sec:unified} is the theorem
\path{coercive_certificate} of \texttt{MovingSofaStability}, and the theorem
\path{gerver_sofa_optimal_unique_stable} of \texttt{MovingSofaExtremal} states
\cref{thm:unified}. The stability library uses the certificate and the theorems
of \texttt{MovingSofaExtremal} at the three places of
\cref{sec:unified-stability}. So \texttt{Solution.lean}, and not only
\texttt{SolutionCoercive.lean}, proves \cref{thm:stability,thm:stability-angle}
through the theorems of \texttt{MovingSofaExtremal}; its proof of
\cref{thm:stability-exponent}\ref{thm:stability-exponent:no} uses neither
optimality nor uniqueness. A third audit,
\path{scripts/AuditCoerciveRoute.lean}, which the continuous integration runs,
follows the proofs through the bodies of all the declarations of the repository,
private and generated ones included. It checks that each of the 743 declarations
of \texttt{MovingSofaExtremal}, \texttt{MovingSofaStability} and
\texttt{SolutionCoercive.lean} uses only the three axioms above. It also checks
that none of them reaches Baek's theorem \baek{Thm.~1.1.1}, the results listed
in the previous paragraph, \texttt{Solution.lean}, or a declaration of the
module \texttt{Main} of \texttt{MovingSofaUniqueness}, of the module
\texttt{Rigidity} (which proves \cref{lem:equality,lem:tangent} and, with
\texttt{Main}, \cref{prop:kernel}) or of the module \texttt{MaximizerRoute} of
the second proof. The modules that prove optimality and uniqueness are
\texttt{Maximizing} of \texttt{MovingSofaUniqueness} and \texttt{Main} of
\texttt{MovingSofaExtremal}. The audit checks that none of their 41
declarations reaches the module \texttt{Margins} of \texttt{MovingSofaStability},
where the stability proof continues after the certificate, or one of the eleven
modules that import it, which include those of the stability theorems and both
solutions; so the argument is not circular. Twenty further checks
confirm that the proofs use what they should; for instance, the classification
of the maximizing caps reaches the certificate and \cref{prop:stab-certificate},
and the stability theorems reach the certificate and the uniqueness theorem of
\texttt{MovingSofaExtremal}. Eight checks confirm that the search finds, in the
older proofs, results that these are known to use, which shows that the search
sees proof bodies. Finally, the audit checks that the fifteen theorems of
\texttt{SolutionCoercive.lean} have the types of those of
\texttt{Solution.lean}.

\subsection{The use of AI}

AI systems wrote the first version of the proof, the formalization and this
text, at the author's request; none of them is an author of this paper.
Appendix~\ref{app:ai} names them, says what each did and what has been checked
by whom.
The argument has not yet been refereed.

\subsection{Open questions}

\emph{Stability constants.} The constants of
\cref{thm:stability,thm:stability-angle} are not explicit; the threshold
$\eps_0$, for instance, comes from a compactness argument
(\cref{prop:stab-entry}). The best constants are not known.
\Cref{thm:stability-exponent} shows that the exponent $1/2$ is optimal for the
Hausdorff distance; whether it is optimal for the area of the symmetric
difference is not known.

\emph{Other corners.} For a hallway whose corner has an angle $\alpha\ne\pi/2$
one can ask for the largest area as a function of $\alpha$ and for the extremal
shapes. Each step of Baek's method has an analogue to prove, and its numerical
inequalities, such as those of \cref{fact:corner}, are specific to the right
angle.

\emph{The ambidextrous sofa.} Romik~\cite{Romik18} found a sofa of area
$1.6449\ldots$ that can turn both left and right around a right-angled corner.
As far as we know, it is not known whether it is optimal. Baek's reduction to
monotone sofas uses one hallway; for a sofa that must turn both ways, the cap,
the niche and the bound $\cQ$ all need two-sided analogues.

\subsection*{Acknowledgments}
This paper builds on the results of~\cite{Baek}. The audit of~\cite{Baek} owes
twelve findings to the notes of another formalization of that
paper~\cite{Cureton}, and was also compared with the blueprint of a
third~\cite{Cao}.
